\section*{Chapitre \ref{ch:g11:deriv} --- Dérivation~: un premier cours}

\begin{solution}{exo:g11:deriv:1}
Pour $f(x) = x^2$ en $a = 2$~:
\[
\frac{(2+h)^2 - 4}{h} = \frac{4h + h^2}{h} = 4 + h
\xrightarrow[h\to0]{} 4,
\]
donc $f'(2) = 4$. Pour $g(x) = x^2 + 3x$ en $a = 1$~: $g(1) = 4$ et
\[
\frac{(1+h)^2 + 3(1+h) - 4}{h} = \frac{5h + h^2}{h} = 5 + h
\xrightarrow[h\to0]{} 5,
\]
donc $g'(1) = 5$.
\end{solution}

\begin{solution}{exo:g11:deriv:2}
$f'(x) = 12x^2 - 4x + 1$.

$g'(x) = \frac{5}{2\sqrt x} - \frac{3}{x^2}$ (sur $\intoo{0}{+\infty}$).

$h'(x) = 2(x^2 - 3) + (2x+1)(2x) = 2x^2 - 6 + 4x^2 + 2x
= 6x^2 + 2x - 6$ (règle du produit).
\end{solution}

\begin{solution}{exo:g11:deriv:3}
$f(x) = x^2 - 3x$, $a = 2$~: $f(2) = -2$, $f'(x) = 2x - 3$ donc
$f'(2) = 1$~; \omterm{def:g11:deriv:tangent}{tangente} $y = -2 + (x - 2) = x - 4$.

$f(x) = \frac1x$, $a = 1$~: $f(1) = 1$, $f'(1) = -1$~; \omterm{def:g11:deriv:tangent}{tangente}
$y = 1 - (x - 1) = 2 - x$.

$f(x) = \sqrt x$, $a = 4$~: $f(4) = 2$, $f'(4) = \frac{1}{2\sqrt4} =
\frac14$~; \omterm{def:g11:deriv:tangent}{tangente} $y = 2 + \frac14(x - 4) = \frac{x}{4} + 1$.
\end{solution}

\begin{solution}{exo:g11:deriv:4}
$f'(x) = \frac{(x-1) - (x+2)}{(x-1)^2} = \frac{-3}{(x-1)^2}$.

$g'(x) = \frac{2x(x^2+1) - x^2 \times 2x}{(x^2+1)^2}
= \frac{2x}{(x^2+1)^2}$.

$h'(x) = -\frac{2x+1}{(x^2+x+1)^2}$ (règle du quotient avec $u = 1$).
\end{solution}

\begin{solution}{exo:g11:deriv:5}
$f'(x) = 3x^2 - 12x + 9 = 3(x^2 - 4x + 3) = 3(x-1)(x-3)$~: positif à
l'extérieur de $\intcc{1}{3}$, négatif à l'intérieur. Donc $f$ est
\omterm{def:g11:func:monotone}{croissante} sur $\intoc{-\infty}{1}$, \omterm{def:g10:functions:variations}{décroissante} sur $\intcc{1}{3}$,
\omterm{def:g11:func:monotone}{croissante} sur $\intco{3}{+\infty}$, avec un \omterm{def:g10:functions:extrema}{maximum} local
$f(1) = 1 - 6 + 9 - 2 = 2$ et un \omterm{def:g10:functions:extrema}{minimum} local
$f(3) = 27 - 54 + 27 - 2 = -2$.
\end{solution}

\begin{solution}{exo:g11:deriv:6}
Sur $\intoo{-1}{+\infty}$~:
\[
f'(x) = \frac{2x(x+1) - (x^2+3)}{(x+1)^2}
= \frac{x^2 + 2x - 3}{(x+1)^2}
= \frac{(x+3)(x-1)}{(x+1)^2}.
\]
Pour $x > -1$, $x + 3$ et $(x+1)^2$ sont tous deux positifs, donc
$f'(x)$ a le signe de $x - 1$~: $f$ est \omterm{def:g10:functions:variations}{décroissante} sur
$\intoc{-1}{1}$ et \omterm{def:g11:func:monotone}{croissante} sur $\intco{1}{+\infty}$, avec \omterm{def:g10:functions:extrema}{minimum}
$f(1) = \frac{4}{2} = 2$.
\end{solution}

\begin{solution}{exo:g11:deriv:7}
$f'(x) = 2x - 4$.

\emph{1.} \omterm{def:g11:deriv:tangent}{Tangente} horizontale~: $f'(x) = 0$ en $x = 2$~; le point est
$(2, f(2)) = (2, 1)$ (le sommet de la \omterm{def:g11:quad:parabola}{parabole}).

\emph{2.} Parallèle à $y = 2x + 1$ signifie pente $2$~: $2x - 4 = 2$,
donc $x = 3$, donnant le point $(3, f(3)) = (3, 2)$.
\end{solution}

\begin{solution}{exo:g11:deriv:8}
Avec largeur $x$ (deux côtés clôturés) et longueur $60 - 2x$~:
$A(x) = x(60 - 2x) = 60x - 2x^2$ sur $\intoo{0}{30}$. Alors
$A'(x) = 60 - 4x$, positif avant $x = 15$ et négatif après~: l'aire est
maximale pour $x = 15$, donnant un enclos $15 \times 30$~m d'aire
$450$~m$^2$. La formule du sommet $\alpha = -\frac{60}{2 \times (-2)} =
15$ du \cref{ch:g11:quad} donne la même réponse sans calcul
différentiel.
\end{solution}

\begin{solution}{exo:g11:deriv:9}
Soit $f(x) = \frac{x+1}{2} - \sqrt x$ sur $\intoo{0}{+\infty}$. Alors
\[
f'(x) = \frac12 - \frac{1}{2\sqrt x} = \frac{\sqrt x - 1}{2\sqrt x},
\]
qui a le signe de $\sqrt x - 1$~: négatif sur $\intoo{0}{1}$, positif
sur $\intoo{1}{+\infty}$. Donc $f$ décroît puis croît, avec \omterm{def:g10:functions:extrema}{minimum}
$f(1) = 1 - 1 = 0$. D'où $f(x) \geq 0$ pour tout $x > 0$,
c'est-à-dire $\sqrt x \leq \frac{x+1}{2}$, avec égalité seulement en
$x = 1$.
\end{solution}

\begin{solution}{exo:g11:deriv:10}
De $\pi r^2 h = 500$, $h = \frac{500}{\pi r^2}$, donc
\[
S(r) = 2\pi r^2 + 2\pi r \cdot \frac{500}{\pi r^2}
= 2\pi r^2 + \frac{1000}{r},
\qquad
S'(r) = 4\pi r - \frac{1000}{r^2}.
\]
$S'(r) = 0$ lorsque $4\pi r^3 = 1000$, c'est-à-dire
$r^3 = \frac{250}{\pi}$, donc $r = \sqrt[3]{250/\pi}$
($\approx 4.3$~cm). Comme $S'(r) = \frac{4\pi r^3 - 1000}{r^2}$ est
négatif avant cette valeur et positif après, c'est un \omterm{def:g10:functions:extrema}{minimum}.
\end{solution}

\begin{solution}{exo:g11:deriv:11}
La \omterm{def:g11:deriv:tangent}{tangente} à $y = x^2$ au point d'\omterm{def:g10:coordgeom:system}{abscisse} $a$ est
$y = a^2 + 2a(x - a) = 2ax - a^2$. Elle passe par $P(1, -3)$ lorsque
$-3 = 2a - a^2$, c'est-à-dire $a^2 - 2a - 3 = 0$, dont les \omterm{def:g11:quad:discriminant}{racines} sont
$a = 3$ et $a = -1$ ($\Delta = 16$). Il y a donc \emph{deux} \omterm{def:g11:deriv:tangent}{tangentes}
par $P$~:
\[
y = 6x - 9 \quad (a = 3),
\qquad
y = -2x - 1 \quad (a = -1).
\]
\end{solution}

\begin{solution}{exo:g11:deriv:12}
D'après l'\cref{ex:g11:deriv:study}, $f$ croît jusqu'au \omterm{def:g10:functions:extrema}{maximum} local
$f(-1) = 3$, décroît jusqu'au \omterm{def:g10:functions:extrema}{minimum} local $f(1) = -1$, puis croît à
nouveau~; pour $x$ grand positif (resp.\ négatif), $x^3$ domine et $f$
prend des valeurs arbitrairement grandes (resp.\ petites). En lisant les
droites horizontales $y = k$ sur le \omterm{def:g10:functions:table}{tableau de variations}~:
\begin{itemize}
  \item $k > 3$ ou $k < -1$~: la droite rencontre la courbe une fois
        --- $1$ solution~;
  \item $k = 3$ ou $k = -1$~: la droite touche un point de retournement
        et croise une autre branche --- $2$ solutions~;
  \item $-1 < k < 3$~: la droite croise les trois branches --- $3$
        solutions.
\end{itemize}
\end{solution}

\begin{solution}{pb:g11:deriv:1}
\textbf{1.} $f'(x) = 6x - 5$~; $g'(x) = 3x^2 - 12$. Par la
définition~: $\frac{(1 + h)^2 - 1}{h} = 2 + h \to 2$~: la \omterm{def:g11:deriv:derivative}{dérivée} de
$x^2$ en $1$ vaut $2$.

\textbf{2.} $g(2) = 8 - 24 = -16$ et $g'(2) = 12 - 12 = 0$~: la
\omterm{def:g11:deriv:tangent}{tangente} est la droite horizontale $y = -16$. Une \omterm{def:g11:deriv:tangent}{tangente}
horizontale signale un point critique --- ici un \omterm{def:g10:functions:extrema}{minimum} local,
comme le confirme le \omterm{def:g10:functions:table}{tableau de variations} de la question 5.

\textbf{3.} \omterm{def:g10:reffunc:affine}{Coefficient directeur} $2a$ en $(a, a^2)$~:
$y = a^2 + 2a(x - a) = 2ax - a^2$.

\textbf{4.} En $x = 0$~: $y = -a^2$~: la \omterm{def:g11:deriv:tangent}{tangente} coupe l'axe à la
profondeur $a^2$ sous l'origine --- exactement aussi bas
\emph{au-dessous} de zéro que $P$ se trouve au-dessus. Recette~:
pour tracer la \omterm{def:g11:deriv:tangent}{tangente} en un point de hauteur $h$, le joindre au
point de l'axe situé à la profondeur $h$ sous l'origine. (Aucune
\omterm{def:g11:deriv:derivative}{dérivée} n'est nécessaire à la planche à dessin.)

\textbf{5.} $g'(x) = 3(x^2 - 4)$~: positive sur
$\intoo{-\infty}{-2}$, négative sur $\intoo{-2}{2}$, positive
au-delà. $g$ croît jusqu'à un \omterm{def:g10:functions:extrema}{maximum} local $g(-2) = 16$, décroît
jusqu'à un \omterm{def:g10:functions:extrema}{minimum} local $g(2) = -16$, puis croît de nouveau.

\textbf{6.} $F\left(0, \frac14\right)$ et $T(0, -a^2)$~:
$FT = \frac14 + a^2$.

\textbf{7.} $FP = a^2 + \frac14 = FT$~: le triangle est isocèle en
$F$, donc les angles à la base sont égaux~:
$\widehat{FTP} = \widehat{FPT}$.

\textbf{8.} Le rayon incident en $P$ est vertical, donc parallèle à
la droite $(TF)$ (l'axe des ordonnées). La \omterm{def:g11:deriv:tangent}{tangente} $(TP)$ coupe ces
deux parallèles, si bien que l'angle entre le rayon et la \omterm{def:g11:deriv:tangent}{tangente}
est égal à $\widehat{FTP}$ (angles alternes-internes) --- lequel
égale $\widehat{FPT}$, l'angle entre $[PF]$ et la \omterm{def:g11:deriv:tangent}{tangente}. Angles
égaux de part et d'autre de la \omterm{def:g11:deriv:tangent}{tangente} en $P$.

\textbf{9.} D'après la loi de la réflexion, le rayon qui arrive
verticalement en $P$ repart selon la droite formant l'angle égal de
l'autre côté --- c'est-à-dire, d'après la question 8, selon $[PF]$~:
par le \omterm{pb:g11:quad:1}{foyer}, quel que soit $a$. En renversant le sens, la lumière
née en $F$ quitte chaque point du miroir parallèlement à l'axe. Les
antennes collectent, les phares éclairent~: démontré.

\textbf{10.} Pour $a = 1$~: \omterm{def:g11:deriv:tangent}{tangente} $y = 2x - 1$, $T(0, -1)$~;
$FP = \sqrt{1 + \left(1 - \frac14\right)^2} =
\sqrt{\frac{25}{16}} = \frac54$~; $FT = \frac14 + 1 = \frac54$~:
isocèle, comme promis.

\textbf{11.} La \omterm{def:g11:deriv:tangent}{tangente} en $x_n$ a pour \omterm{def:g10:algebra:equation}{équation}
$y = f(x_n) + f'(x_n)(x - x_n)$~; elle coupe $y = 0$ en
$x = x_n - \frac{f(x_n)}{f'(x_n)}$.

\textbf{12.} $x_1 = 1 - \frac{-1}{2} = \frac32$~;
$x_2 = \frac32 - \frac{1/4}{3} = \frac{17}{12}$~;
$x_3 = \frac{577}{408}$~: les \omterm{def:g10:numbers:approx}{approximations} de $\sqrt2$ dues à
Héron, c'est-à-dire la machine à \omterm{pb:g10:functions:1}{point fixe} du
\cref{pb:g10:functions:1}. Identité~:
$x - \frac{x^2 - 2}{2x} = \frac{2x^2 - x^2 + 2}{2x}
= \frac{x^2 + 2}{2x} = \frac12\left(x + \frac2x\right)$ --- la
recette de Héron \emph{est} la méthode de Newton appliquée à
$x^2 - 2$, avec deux millénaires d'avance.

\textbf{13.} $x_1 = 5 - \frac{125 - 100}{75} = \frac{14}{3}
\approx 4.6667$~; $x_2 \approx 4.6667 - \frac{1.63}{65.33}
\approx 4.6417$. Calculatrice~: $\sqrt[3]{100} = 4.6416\ldots$ ---
quatre décimales exactes en deux pas.

\textbf{14.} $f'(2) = 0$~: la \omterm{def:g11:deriv:tangent}{tangente} au point de départ est
horizontale (question 2) et ne coupe jamais l'axe --- la formule
divise par zéro. Mise en garde~: lancer Newton à l'écart des points
critiques (et, plus généralement, assez près de la \omterm{def:g11:quad:discriminant}{racine}
cherchée).

\textbf{15.} La dichotomie gagne un chiffre binaire par pas~; Newton
\emph{double} grossièrement le nombre de décimales exactes à chaque
pas ($1 \to 3 \to 6$ ici). Lorsque chaque pas coûte cher --- des
millions de paramètres, une navigation en temps réel --- doubler
l'emporte sur avancer pas à pas, et c'est pourquoi la droite
\omterm{def:g11:deriv:tangent}{tangente} tourne à l'intérieur de tant de silicium.

\textbf{16.} $d^2 = 900 - w^2$, donc
$S(w) = w(900 - w^2) = 900w - w^3$, pour $0 < w < 30$.

\textbf{17.} $S'(w) = 900 - 3w^2 = 0$ en
$w = \sqrt{300} = 10\sqrt3 \approx 17.3$ cm ($S' > 0$ avant,
$< 0$ après~: c'est un \omterm{def:g10:functions:extrema}{maximum}). Alors
$d = \sqrt{600} = 10\sqrt6 \approx 24.5$ cm.

\textbf{18.} $\frac{d^2}{w^2} = \frac{600}{300} = 2$, donc
$d = w\sqrt2$~: la hauteur vaut la largeur multipliée par $\sqrt2$
--- le rapport des charpentiers (et la construction par tiers du
diamètre le livre au compas, sans la moindre algèbre sur le
chantier).

\textbf{19.} Poutre carrée~: $w = d = \sqrt{450} \approx 21.2$ cm,
de rigidité $w^3 = 450^{3/2} \approx 9\,546$. Poutre optimale~:
$S = 10\sqrt3 \times 600 = 6\,000\sqrt3 \approx 10\,392$. Gain~:
environ $8.9\,\%$ de rigidité en plus pour le même tronc --- la
\omterm{def:g11:deriv:derivative}{dérivée} paie la scierie.

\textbf{20.} Géométrie~: les angles égaux de la \omterm{def:g11:deriv:tangent}{tangente} ont fait
d'un paraboloïde un collecteur de lumière. Calcul numérique~:
glisser le long des \omterm{def:g11:deriv:tangent}{tangentes} résout les \omterm{def:g10:algebra:equation}{équations} en doublant les
décimales. Optimisation~: une \omterm{def:g11:deriv:derivative}{dérivée} qui s'annule a trouvé le
rectangle le plus résistant du tronc. Un seul objet, trois métiers
--- et la convexité de l'an prochain dira en outre de quel côté de
chaque \omterm{def:g11:deriv:tangent}{tangente} se tient la courbe.
\end{solution}
